% Generated by docs/generate_user_guide.py. Do not edit directly.
% Source: utilities/descriptions.yaml: preprocessing; preprocessing registry

\subsubsection{Normalization}

\begin{longtable}{@{}>{\raggedright\arraybackslash}p{0.22\textwidth}>{\raggedright\arraybackslash}p{0.20\textwidth}>{\raggedright\arraybackslash}p{0.50\textwidth}@{}}
\toprule
\textbf{Option} & \textbf{Internal key} & \textbf{Description} \\
\midrule
\endfirsthead
\toprule
\textbf{Option} & \textbf{Internal key} & \textbf{Description} \\
\midrule
\endhead
Min-max normalization & min\_max & Converts response values to a 0-1 scale using the observed minimum and maximum within each dataset. With supplied SD/SEM, the range must be finite and positive; errors are divided by that range, which is treated as fixed. \\
Initial-value normalization & initial\_value & Divides each response by the first response value in that dataset when it is nonzero. Supplied SD/SEM is divided by its absolute value, which is treated as fixed. A zero initial response leaves values and errors unchanged. \\
Reverse response (1 - value) & one\_minus & Applies 1 minus each response after any selected normalizations and before logarithmic scaling or interpolation. For data already expressed as fractions, this converts fraction lost to fraction remaining without requiring another normalization. The same selection applies to both calibration profiles and prediction inputs; values are not clipped to 0-1. Supplied SD/SEM is unchanged by this reversal. \\
\bottomrule
\end{longtable}

\subsubsection{Scaling}

\begin{longtable}{@{}>{\raggedright\arraybackslash}p{0.22\textwidth}>{\raggedright\arraybackslash}p{0.20\textwidth}>{\raggedright\arraybackslash}p{0.50\textwidth}@{}}
\toprule
\textbf{Option} & \textbf{Internal key} & \textbf{Description} \\
\midrule
\endfirsthead
\toprule
\textbf{Option} & \textbf{Internal key} & \textbf{Description} \\
\midrule
\endhead
Log10 time scaling & log\_x & Uses a base-10 log scale for time values before fitting or interpolation. All analyzed times must be positive. The browser falls back to unscaled time when calibration times include zero or negative values and reports this in Data Processing. Prediction data must support the applied scaling; CLI inputs with invalid times are rejected. \\
Log10 response scaling & log\_y & Uses a base-10 log scale for response values before fitting or interpolation. This option is unavailable with supplied SD/SEM. Without supplied uncertainty, all processed responses must be positive. The browser can fall back to unscaled response for invalid calibration values and reports this in Data Processing; prediction data must support the applied scaling, and CLI inputs with invalid responses are rejected. \\
\bottomrule
\end{longtable}

\subsubsection{Interpolation}

\begin{longtable}{@{}>{\raggedright\arraybackslash}p{0.22\textwidth}>{\raggedright\arraybackslash}p{0.20\textwidth}>{\raggedright\arraybackslash}p{0.50\textwidth}@{}}
\toprule
\textbf{Option} & \textbf{Internal key} & \textbf{Description} \\
\midrule
\endfirsthead
\toprule
\textbf{Option} & \textbf{Internal key} & \textbf{Description} \\
\midrule
\endhead
No interpolation & none & Uses only exact matches from the observed data. No response-at-time or time-at-response estimates are added, so unmatched positions remain unavailable for Approaches 1 and 3. Approach 2 uses observed points regardless of this setting. \\
Shared-grid interpolation & default & Uses linear interpolation in both directions: response values are estimated at the shared set of observed time points, and times are estimated at the shared set of observed response values. Estimates outside the measured overlap are unavailable. Approach 2 does not use these interpolated points for fitting. \\
Densified interpolation & alternative & Adds intermediate grid points in both time and response-value directions before linear interpolation, producing denser paired mappings for Approaches 1 and 3. The added points are estimates rather than additional experimental observations. Approach 2 does not use them for fitting. \\
\bottomrule
\end{longtable}
